% =============================================================================
\section{Experimental Setup}
\label{sec:setup}
% =============================================================================

\subsection{Datasets}

We evaluate on 31 publicly available datasets spanning five domains: medical,
image, network, tabular, and recommendation or mobility. They range from
several hundred to roughly 200k samples and from 5 to 6169 features. The set
includes the two canonical MIA benchmarks, Purchase100 and Texas100, which
we include specifically to test Finding~2. Datasets larger than 30k samples
are subsampled to 30k to keep the attack-model grid tractable. Most are
loaded through scikit-learn~\cite{pedregosa2011scikit},
OpenML~\cite{vanschoren2014openml}, or the UCI
repository~\cite{dua2019uci}. Table~\ref{tab:all_datasets} lists the
full corpus with sources, and Table~\ref{tab:datasets} details a
representative selection.

\begin{table}[t]
\centering
\caption{The full 31-dataset corpus. OpenML datasets are identified by
         their stable numeric \texttt{data\_id}; the remainder come from
         scikit-learn loaders or their original public releases.}
\label{tab:all_datasets}
\small
\setlength{\tabcolsep}{2.2pt}
\begin{tabular}{lll}
\toprule
\textbf{Dataset} & \textbf{Domain} & \textbf{Source} \\
\midrule
Texas100      & hospital benchmark   & Shokri et al.~\cite{shokri2017membership} \\
Purchase100   & retail benchmark     & Shokri et al.~\cite{shokri2017membership} \\
Adult         & census income        & UCI~\cite{adult1996} \\
COMPAS        & criminal justice     & ProPublica~\cite{compas2016} \\
Pima Diabetes & clinical             & UCI~\cite{smith1988pima} \\
Breast-W      & clinical             & OpenML 15 \\
JM1, KC1      & software defects     & OpenML 1053, 1067 \\
MovieLens-1M  & recommendation       & GroupLens~\cite{movielens} \\
Gowalla       & mobility             & SNAP~\cite{cho2011friendship} \\
MNIST         & image                & OpenML 554~\cite{lecun1998gradient} \\
Fashion-MNIST & image                & OpenML 40996 \\
Digits        & image                & scikit-learn \\
OptDigits     & image                & OpenML 28 \\
PenDigits     & image                & OpenML 32 \\
Letter        & image                & OpenML 6 \\
Segment       & image segmentation   & OpenML 36 \\
SatImage      & remote sensing       & OpenML 182 \\
CovType       & forest cover         & scikit-learn~\cite{blackard1999comparative} \\
German Credit & credit               & OpenML 31~\cite{dua2019uci} \\
Bank Marketing& marketing            & OpenML 1461 \\
Spambase      & e-mail               & OpenML 44 \\
Nomao         & web entities         & OpenML 1486 \\
Electricity   & energy               & OpenML 151 \\
Mushroom      & categorical          & OpenML 24 \\
Magic         & gamma telescope      & OpenML 1120 \\
Phoneme       & speech               & OpenML 1489 \\
HAR           & motion sensors       & OpenML 1478 \\
Gas Drift     & chemical sensors     & OpenML 1476 \\
Vehicle       & silhouettes          & OpenML 54 \\
Ionosphere    & radar                & OpenML 59 \\
\bottomrule
\end{tabular}
\end{table}

The corpus size is a deliberate design choice. Because the dataset is our
unit of statistical analysis, every dataset added to the corpus requires
its own full ground-truth grid of nine attack-model configurations,
including shadow-model training for LiRA, so each point in our regression
costs hours of compute. Thirty-one datasets is roughly an order of
magnitude more than the two to four typically used in MIA evaluations,
and Section~\ref{sec:power} shows empirically why that breadth is not
optional: correlations measured on small corpora can be perfect and
spurious at once.

\begin{table}[t]
\centering
\caption{Representative datasets (of 31). Sample and feature counts are the
         original sizes after preprocessing; datasets larger than 30k
         samples are subsampled to 30k for the attack grid.}
\label{tab:datasets}
\small
\setlength{\tabcolsep}{2pt}
\begin{tabular}{llrrr}
\toprule
\textbf{Dataset} & \textbf{Domain} & \textbf{Samples} & \textbf{Feat.} & \textbf{Cls.} \\
\midrule
Texas100~\cite{shokri2017membership}    & Benchmark/med.  & 67,330 & 6,169 & 100 \\
Purchase100~\cite{shokri2017membership} & Benchmark/retail & 197,324 & 600 & 100 \\
MNIST~\cite{lecun1998gradient}  & Image       & 30,000 & 784 & 10 \\
CovType~\cite{blackard1999comparative} & Forest & 30,000 &  54 &  7 \\
Adult~\cite{adult1996}          & Income      & 48,842 & 14 & 2 \\
German Credit~\cite{dua2019uci} & Credit      &  1,000 & 74 & 2 \\
Spambase~\cite{dua2019uci}      & Spam        &  4,601 & 57 & 2 \\
Gowalla~\cite{cho2011friendship}& Mobility    & 107,092 &  6 & 2 \\
\bottomrule
\end{tabular}
\end{table}

\subsection{Preprocessing}

All datasets pass through one deterministic pipeline. Categorical
features and labels are integer-encoded, every feature is standardized
to zero mean and unit variance, and datasets above 30k samples are
subsampled once, with a fixed seed, before any further processing.
Standardization matters twice over: it is conventional for the models,
and it puts the geometric features of Section~\ref{sec:dpri} on a
common scale so that nearest-neighbour distances are comparable across
datasets. Each dataset is then split 50\%/50\% into members and
non-members by a stratified split with a fixed seed, the standard MIA
evaluation convention; the target model trains on the member half, and
attack AUC is computed over the full member and non-member sets.

\subsection{Ground Truth: Risk(D)}

For each dataset we run three MIA attacks against three model families,
yielding nine AUC scores, and define $\mathrm{Risk}(D)$ as their mean:
\begin{equation}
  \mathrm{Risk}(D) = \frac{1}{|\mathbb{A}||\mathbb{F}|}
  \sum_{\mathcal{A} \in \mathbb{A}} \sum_{f \in \mathbb{F}}
  \mathrm{AUC}(\mathcal{A}, f, D),
  \label{eq:risk_gt}
\end{equation}
where $\mathbb{A} = \{\text{LiRA, shadow, loss-threshold}\}$ and
$\mathbb{F} = \{\text{MLP, XGBoost, RF}\}$.
We use the mean rather than the maximum because the maximum collapses onto
the Random Forest for every dataset (Section~\ref{sec:finding1}), which
saturates near AUC $1.0$ and erases the contrast between datasets. The mean
reflects the typical adversary across realistic and worst-case models alike.

A second metric choice deserves justification. Carlini et
al.~\cite{carlini2022lira} argue that individual attacks should be
judged by true-positive rate at low false-positive rate, because an
attacker who confidently identifies a few members matters more than one
who is weakly right on average. That argument concerns evaluating
\emph{attacks}; our target is a property of \emph{datasets}. We need a
per-dataset scalar that is stable across nine heterogeneous
attack-model configurations and comparable across datasets whose member
sets range from a few hundred to fifteen thousand samples. AUC
integrates over all thresholds, is invariant to the 50/50 class balance
we fix by construction, and remains estimable at our smallest dataset
sizes, where a TPR at $0.1\%$ FPR rests on a handful of samples and is
dominated by noise. Because DPRI predicts a ranking, any monotone risk
functional would do in principle; mean AUC is the one that is stable
enough at $n = 31$ datasets to regress against.

The three attacks span the spectrum from strongest to simplest.
LiRA~\cite{carlini2022lira}, the current state of the art, uses 16
shadow models to compute a per-sample likelihood ratio; the shadow-model
attack~\cite{shokri2017membership} trains 4 shadow models and learns a
binary meta-classifier; and the loss-threshold
attack~\cite{yeom2018privacy} thresholds the model's confidence on the
correct class, which is Bayes-optimal under a Gaussian assumption. The
three model families are an MLP (two hidden layers, 128-64, trained with
Adam~\cite{kingma2015adam} for 100 epochs on GPU),
XGBoost~\cite{chen2016xgboost} (300 trees, depth 6), and a Random
Forest~\cite{breiman2001random} (300 trees, unbounded depth). All models
use a 50\%/50\% member/non-member split, with seeds fixed at 42.
The grid is chosen for coverage rather than exhaustiveness: the attacks
span the simplest baseline to the published state of the art
(Section~\ref{sec:attack_families}), and the model families span the
three regimes that dominate tabular practice, a regularized neural
network, a regularized boosted ensemble, and a deliberately
unregularized memorizer that serves as the adversarial worst case.

\subsection{DPRI Computation}

We compute the five geometric features using scikit-learn and scipy, with
$k=5$ for the $k$-NN features, and append data dimensionality $\log d$. No
model training is required. Silhouette score uses a 5,000-point subsample for
datasets with $n > 5{,}000$.

\subsection{Regression and Validation}

We predict $\mathrm{Risk}(D)$ from the DPRI features under
leave-one-dataset-out cross-validation across all 31 datasets. Features are
rank-transformed within each fold (Section~\ref{sec:dpri}). We report Spearman
rank correlation as the primary metric, together with an unbiased nested
cross-validation that selects the feature subset on the training folds only.
We also report $R^2$, but it is not the appropriate lens here, since the
features are rank-transformed and the prediction target is a ranking.

\subsection{Reproducibility}

All code and data-processing scripts are provided as an anonymized
artifact and will be released publicly upon publication.
All random seeds are fixed, and results are deterministic given the same
hardware.
